\documentclass[11pt]{article}
\usepackage[utf8]{inputenc}
\usepackage[T1]{fontenc}
\usepackage{amsmath}
\usepackage{amssymb}
\usepackage{multicol}
\usepackage{geometry}
\usepackage{tikz}
\usepackage{enumitem}
\usepackage{xcolor}
\usepackage{titlesec}

% Configurações de layout
\geometry{a4paper, left=1cm, right=1cm, top=0.5cm, bottom=1.2cm}
\setlength{\columnseprule}{0.4pt}

% Cores para títulos
\titleformat{\section}{\normalfont\Large\bfseries\color{blue}}{\thesection}{1em}{}
\titleformat{\subsection}{\normalfont\large\bfseries\color{red}}{\thesubsection}{1em}{}

% Título principal - EDITAR CONFORME NECESSÁRIO
\title{\textcolor{red}{[TRIMESTRE]: [TEMA PRINCIPAL]
    \\[SUBTÍTULO OU TÓPICO SECUNDÁRIO]}}
\author{[NOME DO PROFESSOR]}
\date{}

\begin{document}

\maketitle
\vspace{-1cm}

\begin{multicols}{2}

% ============================================================
% SEÇÃO 1: CONCEITOS BÁSICOS
% ============================================================
\section*{Conceitos Básicos}

\subsection*{[PRIMEIRO CONCEITO PRINCIPAL]}
[Definição e explicação do primeiro conceito. Este é um parágrafo explicativo que introduz o tema principal da aula.]

\subsection*{[SEGUNDO CONCEITO PRINCIPAL]}
[Definição e explicação do segundo conceito. Continue com quantos conceitos básicos forem necessários.]

\subsection*{[TERCEIRO CONCEITO PRINCIPAL (se necessário)]}
[Definição e explicação do terceiro conceito.]

\subsection*{Variáveis / Notações}
\begin{itemize}
\item $[variável1]$: [significado da variável 1]
\item $[variável2]$: [significado da variável 2]
\item $[variável3]$: [significado da variável 3]
\item $[variável4]$: [significado da variável 4]
\end{itemize}

% ============================================================
% SEÇÃO 2: PRIMEIRO TÓPICO ESPECÍFICO (ex: Fatorial, Teorema de Pitágoras, etc.)
% ============================================================
\section*{[NOME DO PRIMEIRO TÓPICO]}

\subsection*{Definição}
[Definição clara e completa do primeiro tópico específico.]

\subsection*{Fórmula / Propriedade}
\begin{tikzpicture}
\node[rectangle, fill=blue!10, rounded corners, inner sep=10pt, draw=blue!50, thick] (box) {
$[Fórmula principal em destaque]$
};
\end{tikzpicture}

\subsection*{Exemplo(s)}
[Exemplo resolvido passo a passo mostrando a aplicação do conceito.]
\[
[Desenvolvimento do exemplo]
\]

% ============================================================
% SEÇÃO 3: SEGUNDO TÓPICO ESPECÍFICO (ex: Permutação, Teorema de Tales, etc.)
% ============================================================
\section*{[NOME DO SEGUNDO TÓPICO]}

\subsection*{Definição}
[Definição clara e completa do segundo tópico específico.]

\subsection*{Fórmula / Propriedade}
\begin{tikzpicture}
\node[rectangle, fill=blue!10, rounded corners, inner sep=10pt, draw=blue!50, thick] (box) {
$[Fórmula principal em destaque]$
};
\end{tikzpicture}

\subsection*{Exemplo(s)}
[Exemplo resolvido passo a passo mostrando a aplicação do conceito.]
\[
[Desenvolvimento do exemplo]
\]

% ============================================================
% SEÇÃO 4: TERCEIRO TÓPICO ESPECÍFICO (se necessário)
% ============================================================
\section*{[NOME DO TERCEIRO TÓPICO]}

\subsection*{Definição}
[Definição clara e completa do terceiro tópico específico.]

\subsection*{Fórmula / Propriedade}
\begin{tikzpicture}
\node[rectangle, fill=blue!10, rounded corners, inner sep=10pt, draw=blue!50, thick] (box) {
$[Fórmula principal em destaque]$
};
\end{tikzpicture}

\subsection*{Exemplo(s)}
[Exemplo resolvido passo a passo mostrando a aplicação do conceito.]

% ============================================================
% SEÇÃO 5: COMPARAÇÃO / DIFERENÇAS (opcional)
% ============================================================
\section*{Diferença entre [CONCEITO A] e [CONCEITO B]}

\subsection*{Quando usar [CONCEITO A]}
[A ordem importa / situação específica]. Exemplos:
\begin{itemize}
\item [Exemplo 1 de uso do conceito A]
\item [Exemplo 2 de uso do conceito A]
\item [Exemplo 3 de uso do conceito A]
\end{itemize}

\subsection*{Quando usar [CONCEITO B]}
[A ordem NÃO importa / situação específica]. Exemplos:
\begin{itemize}
\item [Exemplo 1 de uso do conceito B]
\item [Exemplo 2 de uso do conceito B]
\item [Exemplo 3 de uso do conceito B]
\end{itemize}

% ============================================================
% SEÇÃO 6: EXERCÍCIOS BÁSICOS
% ============================================================
\section*{Exercícios Básicos - [PRIMEIRO TÓPICO]}

\begin{enumerate}[leftmargin=*]

\item [Enunciado do exercício 1]

\item [Enunciado do exercício 2]

\item [Enunciado do exercício 3]

\item [Enunciado do exercício 4]

\end{enumerate}

% ============================================================
% SEÇÃO 7: EXERCÍCIOS - SEGUNDO TÓPICO
% ============================================================
\section*{Exercícios Básicos - [SEGUNDO TÓPICO]}

\begin{enumerate}[leftmargin=*, start=5]

\item [Enunciado do exercício 5]

\begin{enumerate}[label=\alph*)]
\item [subitem a]
\item [subitem b]
\item [subitem c]
\end{enumerate}

\item [Enunciado do exercício 6]

\item [Enunciado do exercício 7]

\item [Enunciado do exercício 8]

\end{enumerate}

% ============================================================
% SEÇÃO 8: EXERCÍCIOS - TERCEIRO TÓPICO (se necessário)
% ============================================================
\section*{Exercícios Básicos - [TERCEIRO TÓPICO]}

\begin{enumerate}[leftmargin=*, start=9]

\item [Enunciado do exercício 9]

\begin{enumerate}[label=\alph*)]
\item [subitem a]
\item [subitem b]
\item [subitem c]
\end{enumerate}

\item [Enunciado do exercício 10]

\item [Enunciado do exercício 11]

\end{enumerate}

\end{multicols}

\end{document}
